\begin{center}
    {\large \textbf{Appendice}}
\end{center}

\noindent \newline \newline

Descrizione delle modalità adottate durante il viaggio. 
\newline
\newline

\textbf{Modalità Marines} — anche detta modalità "Sbrigati e aspetta": tutto viene fatto velocemente ma con attenzione, cercando di arrivare il prima possibile alla prossima meta. Perfetta per gli spostamenti più lunghi e specialmente negli aeroporti, dove si corre come se si avessero i razzi sotto i piedi, ma poi si resta inchiodati in fila alla sicurezza.
\newline

\textbf{Modalità Bodyguard} — ideale nei mercatini locali più pericolosi. Occorre vestirsi come un ex-militare americano dei film e proteggere il bersaglio — alias moglie — dalle potenziali minacce circostanti, senza mai toglierle gli occhi di dosso.
\newline

\textbf{Modalità Fantozzi} — quando si è troppo stanchi e si mette la sveglia all'ultimo minuto per recuperare sonno, rischiando di fallire tutti gli appuntamenti successivi, mentre la giornata si trasforma in una serie di disastri inevitabili.
\newline

\textbf{Modalità National Geographic} — perfetta quando si fotografano gli animali. Bisogna pensare che le nostre foto potrebbero finire su un'edizione del National Geographic: più foto si fanno, più si aumentano le probabilità di farne di essere pubblicati.
\newline

\textbf{Modalità Struzzo} — quando si cercano le tartarughe tra le proprie dita dei piedi in mare. Gambe dritte, maschera da sub e si immerge in acqua solo la testa.
\newline

\textbf{Modalità Antiscam} — a costo di essere rompiballe, si fanno tutte le domande necessarie per conoscere a fondo dettagli, prezzi e orari, trasformandosi nel cliente più fastidioso di tutti i tempi pur di evitare truffe.
\newline